%!TeX root = principal.tex
%!TeX encoding = utf-8
Capsule networks were proposed with a promise of equivariance: by explicitly representing the pose of entities, they would be more robust to transformations absent from training and would need less data. Since then, data augmentation has given convolutional and attention-based networks much of that tolerance. This work tests whether the capsule advantage still exists. We compare Efficient-CapsNet, ResNet-18 and DeiT-Tiny, trained from scratch with identical hyperparameters, splits and seeds, on CIFAR-10 and on a binary version of D-Fire, a fire and smoke image dataset. The full factorial grid has 162 training runs, crossing three augmentation regimes, three data fractions and three seeds, and every model is evaluated under large rotations and under a composition of unseen corruptions. The results do not support the promises. Under rotation, the capsule model's accuracy retention does not exceed ResNet-18's on either dataset, and on CIFAR-10 ResNet-18 has higher absolute accuracy in every augmented regime. The capsule model's smaller accuracy drop when data is reduced is about one percentage point and is not significant. Augmentation widens the gap instead of closing it: on CIFAR-10, the capsule deficit on clean data grows from 2.5 to 13.2 percentage points between the no-augmentation and strong regimes. Efficient-CapsNet stands out elsewhere: it is the most stable architecture across seeds, its confidence identifies its own errors well on clean data, and on D-Fire it comes within a few points of the best baseline with a feature extractor tens of times smaller. The work also shows that the retention ratio, a common robustness metric, favours weak models, and that compositions of many corruptions drive all models to chance level.